\documentclass[11pt]{article}
% Shared preamble for per-Python-file documentation (input via \input{...}).
\usepackage[margin=1in]{geometry}
\usepackage[T1]{fontenc}
\usepackage[utf8]{inputenc}
\usepackage{lmodern}
\usepackage{hyperref}
\usepackage{xcolor}
\usepackage{listings}
\usepackage{listingsutf8}
\lstset{
  basicstyle=\ttfamily\small,
  columns=fullflexible,
  breaklines=true,
  upquote=true,
  inputencoding=utf8,
  extendedchars=true,
  frame=single,
  framerule=0.2pt,
  tabsize=2
}


\title{\texttt{run\_d2q.py} (Q\&A)}
\author{}
\date{}

\begin{document}
\maketitle

\paragraph{Q: What is the purpose of \texttt{run\_d2q.py}?}
\textbf{A:} It trains and evaluates the D2Q baseline. The script loads precomputed per-duration-bucket play-time quantiles, trains a model to predict a normalized quantile position, then maps that prediction back to a watch-time value for evaluation.

\paragraph{Q: What extra data does D2Q require beyond \lstinline|kuairec_data.pkl|?}
\textbf{A:} It expects \lstinline|d2q_duration_bucket_playtime_quantiles.csv| (created by preprocessing) and loads it into a \lstinline|bucket_quantiles| dictionary, which defines how values map to and from quantile space within each duration bucket.

\paragraph{Q: How are labels transformed for training?}
\textbf{A:} Each label value is mapped into a bucket-specific quantile index using \lstinline|from_value_to_quantile|, then normalized into \([0,1]\) by \lstinline|label_norm| before computing an MSE loss against the model output.

\paragraph{Q: How are predictions transformed for evaluation?}
\textbf{A:} The model outputs a scalar quantile-like value; the script uses \lstinline|from_quantile_to_value| (bucket-specific interpolation) to map it back to an estimated watch-time value.

\paragraph{Q: Code example: mapping predicted quantile back to value?}
\textbf{A:}
\begin{lstlisting}[language=Python]
mapped_y_value = from_quantile_to_value(buckets_quantiles, idx.item(), quantile)
\end{lstlisting}

\paragraph{Q: Full source code?}
\textbf{A:}
\lstinputlisting[language=Python]{/workspace/run_d2q.py}

\end{document}
